%%
%% manued-e.tex : the manual of manued.el in English
%% Copyright (C) 1998-2012 Hitoshi Yamauchi ���� ��
%%
%% latex2html -html_version 4.0 -init_file dot.latex2html-init-eng manued-e.tex
%%
\documentclass{article}
\usepackage{html}
\usepackage{makeidx}
\usepackage{graphics}
% \usepackage{graphicx}
\usepackage{color}
\makeindex

\title{Manued Reference Manual \\
  \includegraphics{manuedlogo}}
\author{Hitoshi Yamauchi }


\begin{document}
\maketitle

\begin{rawhtml}
<!--#config timefmt="%Y/%m/%d (%A)" -->
\end{rawhtml}

\begin{abstract}
 Manued is a proofreading method of manuscripts, proposed by Ikuo
 Takeuchi in No.39 Programming symposium, 1998. Although, Manued is a
 media independent proofreading method, it is especially effective for
 exchanging electronic text via E-mail. \texttt{manued.el} is an emacs
 lisp implementation to realize this method on emacs.
\end{abstract}
\tableofcontents
\newpage
\begin{rawhtml}
<hr>
\end{rawhtml}

\section{Introduction}
\subsection{Brief description of Manued}

Oddly, even in this digital world, revising electric manuscripts is not
easy as revising them on a piece of paper with red ink. Especially, in
an email environment, it is hard to revise a manuscript since the kinds
of usable characters are limited, or your email client might
automatically change the text arrangement. However, many people want to
exchange manuscripts by email. Conventionally, many people may use a
method for revising an electric manuscript like below.
\begin{verbatim}
    The Three Laws of Robotics by Azimuth.
                                  ^^^^^^^
             This may be a mistake for Asimov. And the number of
             laws was changed from three to four, wasn't it?

    1. A robot may not injure a robot, nor through inaction
                                ^^^^^ This must be "human being"

       allow a human being to come to harm.
\end{verbatim}

This conventional method indicated correcting position with a symbol
`\^' and proofreading plans with a text. This method always needs much
input work and includes some kind of ambiguous, because the proofreading
plans are directed by natural language.  The main problem of this
conventional method comes from this ambiguity. Even if the manuscript
has already been digial data, this ambiguity makes automatic extraction
of the revised document difficult. Furthermore, information of
correcting position with `\^' sometimes may not be kept depends on the
mail client.  For example, someone reads a manuscript with proportional
font and the other doesn't. In other words, there is the case that it is
difficult to specify a correcting point since each person has own
environment to read an email including a manuscript.

Manued method solves these problems. The input for designation of
proofreading is easy, and it can define exact designation of the
correcting region and how to change the text.  Therefore, there is no
failure the correcting position. There is no mis-recognition of the
proofreading text that is a comment or a substitution. In particular,
the manued method is independence from the reading environment. We
should only assume the manuscript is written as a text file.

Manued has two main notation of proofreading rule. One is ``delete and
insert'' rule. This is shown with notation of [original text
(delete) /revised text (insert) ; comments ]. Therefore, above
text will be changed as follows.

\begin{quotation}
  The Three Laws of Robotics by \verb|[|Azimuth/Asimov; This may be a
  mistake with Asimov. And the number of laws was changed from three
  to four, isn't it?\verb|]|.  1. A robot may not injure a
  \verb|[|robot/human being\verb|]|, nor through inaction allow a
  human being to come to harm.
\end{quotation}

Of course, we can use this method with a piece of paper and a
pen. However, we assume only an editor program edits the text file. I
implemented ``manued-minor-mode'' on the Emacs editor to make this
method more useful. In manued-minor-mode, we can use following
functions.
\begin{itemize}
 \item Extracting the original document or the revised one from one
       manued document automatically
 \item Searching the proofreading position
 \item Assistance for inserting manued commands
\end{itemize}

Since manued method doesn't care the writing environment, an input
assisting tools can be implemented on other editors. Just only because I
always use the emacs, I implemented the tool on there. So, please let me
know if you implemented an assistance tool for manued method on other
environment. Examples of demanded environments are as follows.
\begin{itemize}
\item vi
\item Microsoft Word
\item Filter commands for extracting revised document, etc.
\end{itemize}

\subsection{Information of this document}

\begin{itemize}
\item Copyright (C) 1997-2012 Hitoshi Yamauchi
\end{itemize}

\newpage
\begin{rawhtml}
<hr>
\end{rawhtml}

\section{Quick Start of Manued}
% <label id="QuickStart">

\subsection{About Manued}

Manued is a method for proofreading of manuscripts, proposed by Ikuo
Takeuchi in No.39 Programming symposium, Japan, 1998. The aim of this
method is to achieve easy proofreading when people exchange a manuscript
via email.  I present \texttt{manued.el} as one of the implementations
to assist the method.

\subsection{Installation}

\texttt{manued.el} depends on the Emacs version 20.7 or
later. Installation procedure of \texttt{manued.el} is as following.

\begin{enumerate}
 \item Put the \texttt{manued.el} file on your emacs `load-path'.
 \item Add the next line to your .emacs file.
\begin{verbatim}
(autoload 'manued-minor-mode "manued" "manuscript editing mode" t)
\end{verbatim}
\end{enumerate}


\subsection{Starting Manued}

To try out the \texttt{manued.el}, you can read the tutorial
file(\texttt{manued.tut}) or directly enter to the
\texttt{manued-minor-mode}.

\begin{itemize}
\item Read the tutorial \\ %
  You can find the \texttt{manued-e.tut} file in the
  manued package. Please read the \texttt{manued-e.tut} file in your
  emacs buffer.
\item Directly enter to the \texttt{manued-minor-mode} \\ %
  \begin{enumerate}
  \item To enter \texttt{manued-minor-mode}, type ``M-x
    manued-minor-mode''.
  \item Type some texts in the buffer.
  \item To insert a manued swap command, put a mark and type M-x
    manued-insert-swap-command (or type C-c C-m C-s).
  \item To extract the revised document in another buffer, Type
    manued-show-newer-in-manued-buffer.
  \item Others: select from a menu item. Type M-n or M-p, this is
    the search command for manued command.
  \end{enumerate}
\end{itemize}

\newpage
\begin{rawhtml}
<hr>
\end{rawhtml}

\section{What is Manued?}
\subsection{Outline of Manued}

Manued is a proofreading method of manuscripts, proposed by
\htmladdnormallink{Ikuo Takeuchi}{http://ne.cs.uec.ac.jp/tak.html} in
No.39 Programming symposium~\cite{bib:Takeuchi1998}, 1998. Manued stands
for MANUscript EDiting Language (and there is another origin. Takeuchi
is interested in `Nue'. This is an imaginary creature of
Japan. Ma-nue-dou has special meaning in Japanese). The aim of this
method is to achieve easy proofreading with exchanging a manuscript via
email.(\htmladdnormallink{Here is original text in Japanese on the
Web.}{./manuscripteditting.txt.gz})

When we revise a manuscript sent by email, there are some problems as
following.
\begin{itemize}
 \item We cannot use the special symbols like hand written characters
       for the proofreading a manuscript. The email usable-character-set
       is limited.

 \item Using colors are cumbersome.

 \item Some mail reader doesn't keep two dimensional information of the
      document. When the author and the editor use different editor,
      there is no guarantee that the positioning information of the
      document is the same. Originally email protocol (RFC 822) does not
      have any rules about character width, fonts and how to display
      breaking lines. Therefore, these situations can usually occur.

 \item We don't want to depend on the line feed position. To send exact
       correction information, we would like to avoid the effect of the
       editor software's line feed processing algorithm for a long line,
       which depends on each editor program.
\end{itemize}

To solve these problems, we used to exchanging printed paper including
proofreading information with red ink via snail-mail. However, this
method takes too long in this digital era. Using facsimile is another
solution since it has acceptable transmission speed, but receiver must
inputs corrections by reading pieces of paper with his/her eyes. So,
some errors will be mixed to a manuscript by human error, and input
costs for human being is high. There is the OCR technology for
recognizing hand written characters, but why we need to use them if we
have already typed in the email. In the end, it seems not a good idea to
change the media of the manuscript from electric stuffs to printed
papers.

\cite{bib:Takeuchi1998} pointed out that proofreading for an electric
text should have the following features.
\begin{itemize}
 \item Both revised document and original one are extracted
       automatically. (required)
 \item High readability for human being.
 \item Correction points should be easily found.
 \item We can put some comments for revised region.
 \item It should be easy to put revision information to a manuscript.
 \item Repeated correction directive should be abstracted.
\end{itemize}

Manued is proposed to satisfy these features and this stands for
MANUscript EDiting language. The author presents an implementation to
support proofreading with this language.


\subsection{Examples of Manued document.}
% <label id="reidai">

Let me show you some simple examples.

\begin{latexonly}
  \begin{itemize}
   \item He was {\color{red}\textbf{\underline{killed}}} at programming.
   \item When ``Norwegian {\color{red}\textbf{\underline{woods}}}'' is
         translated to Japanese, it is always
         {\color{red}\textbf{\underline{interpretented}}} as ``a forest in
         Norway''. However, {\color{red}\textbf{\underline{1}}} of my
         friends says that the {\color{red}\textbf{\underline{light}}}
         meaning is a furniture made in Norway. Please tell me which is
         right?
   \item Why do people say ``chapter
         {\color{red}\textbf{\underline{first}}}''?
  \end{itemize}
\end{latexonly}
\begin{rawhtml}
  <ul>
  <li>He was <b><u><font color="#ff0000">killed</font></u></b>
  at programming.
  <li>When ``Norwegian <b><u><font color="#ff0000">woods
  </font></u></b>'' is
  translated to Japanese, it is always
  <b><u><font color="#ff0000">interpretented</font></u></b>
  as ``a forest in Norway''. However,
  <b><u><font color="#ff0000">1</font></u></b> of my friends says that the
  <b><u><font color="#ff0000">light</font></u></b>
  meaning is a furniture made in Norway. Please tell me which is
  right?
  <li>Why do people say ``chapter <b><u><font
  color="#ff0000">first</font></u></b>''?
  </ul>
\end{rawhtml}


I emphasized the point which I would like to correct in this
text. However, some browsers may not express them. I want to correct
above text as follows.

\begin{latexonly}
  \begin{itemize}
  \item He was {\color{red}\textbf{\underline{skilled}}} at
    programming.

  \item When ``Norwegian {\color{red}\textbf{\underline{wood}}}'' is
    translated to Japanese, it is always
    {\color{red}\textbf{\underline{interpreted}}} as ``a forest in
    Norway''. However, {\color{red}\textbf{\underline{one}}} of my
    friend says that the {\color{red}\textbf{\underline{right}}}
    meaning is a furniture made in Norway. Please tell me which is
    right?
  \item Why do people say ``chapter
    {\color{red}\textbf{\underline{one}}}''?
  \end{itemize}
\end{latexonly}
\begin{rawhtml}
  <ul>
  <li>He was <b><u><font color="#ff0000">skilled</font></u></b> at programming.
  <li>When ``Norwegian <b><u><font color="#ff0000">wood</font></u></b>'' is
  translated to Japanese, it is always
  <b><u><font color="#ff0000">interpreted</font></u></b>
  as ``a forest in Norway''. However,
  <b><u><font color="#ff0000">one</font></u></b>
  of my friend says that the
  <b><u><font color="#ff0000">right</font></u></b>
  meaning is a furniture made in Norway. Please tell me which is
  right?
  <li>Why do people say ``chapter
  <b><u><font color="#ff0000">one</font></u></b>''?
  </ul>
\end{rawhtml}

By the conventional method, we correct the text as the following.

\begin{verbatim}
He was killed at programming.
       ^ skilled. This small missing "s" is fatal!


When ``Norwegian woods'' is translated to Japanese, it is always
                     ^ delete s. I love the Beatles.

interpretented as ``a forest in Norway''. However, 1 of my friends
^ interpreted.                                     ^ Of course one is one,
                                                     but write `one'.

says that the light meaning is a furniture made in Norway.
              ^ right.

Please tell me which is right?


Why do people say ``chapter first''?
                            ^ one. (I have also same question but
                                    one is right.)
\end{verbatim}


To display this conventional descriptions for browsers, I use special
tags to preserve two dimensional information like verbatim (for LaTeX)
or pre (for html). However, in email environment (RFC 822), there is no
such special directives and this conventional descriptions using `\^'
symbol may corrupted. For example, the same information will be shown in
the next paragraph without the tags which preserves two dimensional
information. The representation of above text will be varied with each
environment, since usual browser ignores successive white spaces and
line feeds. Some people can understand next documents, but the others
see meaningless texts. It will be clear at a glance that this is not
beneficial for the people who can not read the following text.

He was killed at programming.
       \^ skilled. This small missing "s" is fatal!


When ``Norwegian woods'' is translated to Japanese, it is always
                    \^ delete s. I love the Beatles.

interpretented as ``a forest in Norway''. However, 1 of my friends
\^ interpreted.                                   \^ Of course one is one,
                                                     but write `one'.

says that the light meaning is a furniture made in Norway.
              \^ right.

Please tell me which is right?


Why do people say ``chapter first''?
                            \^ one. (I have also same question but
                                     one is right.)


In addition, even if an environment can preserve the two dimensional
information, it is hard to distinguish among correction information,
only the comment and original text.  Moreover, first you corrected the
latter part of your manuscript and second you correct former part of the
manuscript, latter corrected parts are sometimes shifted. In such
situation, you must maintain the manuscript for keeping consistency of
two dimensional information of revised position. This is not easy.  I
stressed also that cut and paste manipulation may be harmful for the
conventional method.  Proofreading with Manued is independent from two
dimensional information and can resolve the problems of conventional
method using '\^' symbols.

In Manued method, there are only two replacement rules shown as the
following.

\begin{itemize}
\item \verb-[A/B]- : This means `delete A and insert B'.
  The variations are
  \begin{enumerate}
  \item \verb-[A/]- : Delete A (delete A and insert nothing).
  \item \verb-[/A]- : Insert A (delete nothing and insert A).
  \end{enumerate}
\item \verb-[A|B|C]- : Swap A with C, so result will be `CBA'.  \\ %
  The variation is
  \begin{enumerate}
  \item \verb-[A||C]- : Swap A with C.
  \end{enumerate}
\end{itemize}

According to above rules, the example texts will be presented as
the following.

\begin{latexonly}
  \begin{itemize}
  \item He was {\color{red}\textbf{\underline{[/s; This small missing
          "s" is fatal!}]}}killed at programming.
  \item When ``Norwegian wood {\color{red}\textbf{\underline{[s/; I
          love the Beatles.}]}}''  is translated to Japanese, it is
    always
    {\color{red}\textbf{\underline{[interpretented/interpreted]}}} as
    ``a forest in Norway''. However,
    {\color{red}\textbf{\underline{[1/one; Of course one is one, but
          write `one'.]}}}  of my friends says that the
    {\color{red}\textbf{\underline{[light/right]}}} meaning is a
    furniture made in Norway. Please tell me which is right?
  \item Why do people say ``chapter \\
    {\color{red}\textbf{\underline{[first/one; I have also same
          question but one is right.]}}}''?
  \end{itemize}
\end{latexonly}
\begin{rawhtml}
  <ul>
  <li>He was <b><u><font color="#ff0000">[/s; This small missing
  "s" is fatal!]</font></u></b>killed at programming.

  <li>When ``Norwegian wood
  <b><u><font color="#ff0000">[s/; I love the Beatles.</font></u></b>''
  is translated to Japanese, it is always
  <b><u><font color="#ff0000">[interpretented/interpreted]</font></u></b>
  as ``a forest in Norway''. However,
  <b><u><font color="#ff0000">[1/one; Of course one is one, but
        write `one'.]</font></u></b>
  of my friends says that the
  <b><u><font color="#ff0000">[light/right]</font></u></b>
  meaning is a furniture made in Norway. Please tell me which is
  right?

  <li>Why do people say ``chapter
  <b><u><font color="#ff0000">[first/one; I have
  also same question but one is right.]</font></u></b>''?
  </ul>
\end{rawhtml}

You can see the method is stronger than the conventional method. I think
you see how manued works. If there is a tool which processes these
rules, manued method will be more powerful and useful. Therefore, I
implemented \texttt{manued.el}, this program supports manued
method. \texttt{manued.el} has the following functions : highlighting
correction regions, searching manued directives, inserting manued
directives, extracting revised document from the original one.

\newpage
\begin{rawhtml}
<hr>
\end{rawhtml}

\section{Syntax of Manued}
\subsection{Manued command and Manued definition commands of
  correction symbols}

In Manued, there are ``definition commands of correction symbols'' and
``manued commands''.  ``Definition commands of correction symbols'' are
to avoid dependences of different environments and ``manued commands''
is to indicate proofreading points. We also call definition commands of
correction symbols as def-command, because these commands are always
started with a string `def'. Def-commands define proofreading symbols in a
manuscript. To correspond to next circumstances, def-commands is defined
by Manued.
\begin{itemize}
 \item Since usable characters are limited in email environment, the
       characters indicating correction points with Manued will be
       varied according to a document.

 \item It is almost same with above. If fixed correction characters are
       used for indicating proofreading point by manued, some
       environment forces to use too many escape commands. This will
       degrade readability of the manued document. We would like to
       avoid these situations using def-commands.
\end{itemize}

Def-command is always put at the beginning of the document and must be
put at the beginning of line. When it is not put at the beginning of
line, Manued could not recognize them in my implementation. The names of
definition commands of correction symbols are `def' plus each manued
command name. For example, `delete and insert' command is defined by
`defdelete' def-command.

When you want to set delete command symbol as `*', you should insert
next line in your document.

\begin{verbatim}
        defdelete *
\end{verbatim}
Then correction region will be written as next.
\begin{verbatim}
        [original text*revised text]
\end{verbatim}
This means delete `original text' and insert `revised text'.

When your document has no def-command, default characters are used
that suggested in \cite{bib:Takeuchi1998} by Ikuo Takeuchi. I will
also show the default def-commands in the following section.

\subsection{Definition commands of correction symbols for
  \texttt{manued.el} ver.0.9}
\label{sec:manueddefcommand}

This section explains the def-commands.

Def-commands always appears in a proofreaded document as the following.
\begin{verbatim}
        defparentheses  [ ]
        defdelete       /
        defswap         |
        defcomment      ;
        defescape
        deforder        newer-last
        defversion      0.9.5
\end{verbatim}
The details of the each line are:
\begin{itemize}
\item defparentheses\\ %
      Defines the start string and end string of manued
      command.  (default value `\texttt{[ ]}')
\item defdelete\\ %
      Defines delete and insert string. In default, delete
      the former part of this symbol and insert the latter part of this
      symbol. For example, \texttt{[A/B]} means delete A and insert
      B. You can exchange the delete part and insert part, see also
      deforder.  (default value `\texttt{/}')
\item defswap\\ %
      Defines swap command. For example,
      \verb-[A|B|C]- means A and C will be swapped and results is
      CBA.  (default value `\verb-|-')
\item defescape\\ %
      Defines manued escape command. For example, [This is
      \verb|~|[foo\verb|~|]/These are \verb|~|[bars\verb|~|]] has
      delete ``This is [foo]'' and insert ``These are [bars]''.  However,
      too many escapes and is too complicated to read. To avoid this
      many escape situation, you can re-define manued command strings by
      using defparentheses.  For example, write ``defparentheses \{
      \}'', we can rewrite below to \{This is [foo]/These are [bars]\}.
      (default value `\verb|~|') \\
      Here, you must pay attention that
      escaped symbol is starting manued command strings only when it is
      not in manued command. This is avoid complicated situation like
      below.
\begin{verbatim}
        The ratio $a_n/b_n$ must be considered here.
        ~[\lim_{n \to \infty} \frac{a_n}{b_n} \]
\end{verbatim}
      If any manued command symbol is processed at any place, all manued
      command in original text must be escaped. In the above example,
      '/' of $a_n/b_n$ and ] is corresponding to manued command
      symbol. But these symbols are not in the manued command. It is too
      cumbersome to escape all the manued command in the
      text. Therefore, `starting manued region strings' is the only
      escape target (in here, `[') in the plain text in my
      implementation.
 \item defcomment\\ %
       Defines comment of manued command. Latter part of this symbol
       in manued command will be ignored.  (default value `\texttt{;}')
 \item deforder \label{item:deforder}\\ %
       This defines newer part of delete command. In defdelete command of
       the manued method, the deleted part is called as the `older' part,
       and the inserted part is called as the `newer' part. The list in
       delete command is ordered as [first/last]. You can decide which is
       delete part.  So, these (newer older) and (first last) combinations
       will be next 4 types.
       \begin{itemize}
        \item newer-first
        \item newer-last
        \item older-first
        \item older-last
       \end{itemize}
       This deforder of definition commands of correction symbols is only
       define attribute of delete command, and doesn't define any manued
       command symbol.  (default value `\texttt{newer-last}')
 \item defversion\\ %
       Shows the version information for \texttt{manued.el}.  This is
       useful information for the \texttt{manued.el} program.  (default
       value `(manued-version-num)')
\end{itemize}

Definition commands should be located at the top of a manuscript and
each def-command must be located at the beginning of line. If a
manuscript does not include these def-commands, the default manued
symbols will be used. \texttt{manued.el} checks consistency of
def-commands and warns to the user when it is not consistent, for
example, if defparentheses has only one argument, it will be warned.

\subsection{manued commands for \texttt{manued.el} version 0.9 or later}

This section explains manued commands with default manued symbols.

\begin{itemize}
 \item \verb-[A/B]-\\ %
       Delete A, insert B The variations are
       \begin{enumerate}
        \item \verb-[A/]- : Delete A. (delete A, insert null string.)
        \item \verb-[/A]- : Insert A. (delete null string, insert A.)
       \end{enumerate}

 \item \verb-[A|B|C]-\\ %
       Swap A with C. This will be CBA.  \\ %
       The variation is
       \begin{enumerate}
        \item \verb-[A||C]- : Swap A with C. This will be CA.
       \end{enumerate}

 \item \verb-[; comment ]- \\ %
       Set a comment range from `;' to `]'.

 \item Note \\
       Sometimes you might want to write the manued command string
       itself. In that case, you can just escape the command. For
       example, \verb|~[| will be \verb|[|.
\end{itemize}

Nesting of these commands is possible, and process is done from the
inside to outside. However, non interpretable list of manued command is
an error.  For example, it is an error that there is no matching
defparenthesis (e.g., not closing or closing without opening), or it is
also an error like \verb-[A|B/C]- manued sequence.

\subsection{Pretty print (after ver.0.9.5)}

When you extract the new document with manued, do you want to modify the
extracted part? For example, when you use a markup language and you want
to emphasize the proofreading part with the language.  If yes, there is
the way to modify the extracted part.

First, you type
\begin{verbatim}
  M-x manued-set-pretty-print-on
\end{verbatim}
and, you apply some Extracting document commands. I will describe this
with the following example.

Example 1: Extract from \verb-[A/B;C]- to
\verb-\colortext{blue}{A}\colortext{red}{B}\colortext{green}{C}-.
Namely, when you want to add some colors at revised parts with TeX.
\begin{verbatim}
(setq manued-pretty-print-format-delete-list
  '("\\colortext{blue}{%s}" "\\colortext{red}{%s}" "\\colortext{green}{%s}"))
\end{verbatim}
This should be included in your \texttt{.emacs} or \texttt{.manued.el}.
The elements of this list correspond with each part of delete command
like next:
\begin{verbatim}
(setq manued-pretty-print-format-delete-list
  '("First Element Format" "Second Element Format" "Comment Element Format"))
\end{verbatim}
Each string element will be an argument of the format function of the Emacs.

\subsubsection{Examples of pretty print}

Ex. 2: In \verb-[A/B;C]-, when you want to output B part with
\texttt{Large} strings in TeX way and want to delete part A and comment
part.
\begin{verbatim}
 (setq manued-pretty-print-format-delete-list
   '("" "{\\Large %s}" ""))
\end{verbatim}

Ex. 3: If you want to change the style of manued expression like
from \verb-[A/B;C]- to \verb-{A*B@C}-.

\begin{verbatim}
(setq manued-pretty-print-format-delete-list
  '("{%s*" "%s" "@%s}"))
\end{verbatim}

In manued, the comment part is an optional (not needed).  Therefore,
when there is no comment in manued expression, the comment part format
string is not outputted by default.  This behavior is controlled by the
variable \verb|manued-pretty-print-null-comment-out-p|.  Here, we
describe only about delete command, but also the way to control swap
command with pretty print mode. See \ref{sec:PrettyPrintCustomize} in
details.


\subsection{Examples of manued documents}
\subsubsection{Example 1}

Since defparenthesis can indicate any strings without white space, next
example is valid.
\begin{verbatim}
    defparentheses ^_^[  ]^_^

    Manued^_^[,/]^_^ stands for ``^_^[Tao of manued/Manuscript Editing
    Language; In Japanese, we call this Manuedou and dou means Tao.]^_^''.
\end{verbatim}

\begin{itemize}
\item Original version :\\ %
      Manued, stands for ``Tao of manued''
\item Revised version : \\ %
      Manued stands for ``Manuscript Editing Language''.
\end{itemize}

\subsubsection{Example 2}
When you have an environment that can use multi-byte characters, you may
use some def-commands with multi-byte characters. This document of
Japanese version shows such examples. But this document does not assume
such environment, I omit the example multi-byte def-command.

\subsubsection{Example 3}

The following text looks like a SGML, HTML, XML.

\begin{verbatim}
    defparentheses      <manued>  </manued>
    defdelete   /

    ``Don't worry, Miss Sakuara. The manuscript still is.
      <manued>/At here.</manued>''

    ``Give me the supreme <manued>sing/song.</manued>''

    `` <manued>Reverse/Return</manued> to your original figure.''
\end{verbatim}
%% `` <manued>'s first space is needed for latex2html

\begin{itemize}
\item Original version :\\ %
    ``Don't worry, Miss Sakuara. The manuscript still is.''
    ``Give me the supreme sing''
    ``Reverse to your original figure.''
\item Revised version : \\ %
    ``Don't worry, Miss Sakuara. The manuscript still is. At here.''
    ``Give me the supreme song.
    ``Return to your original figure.''
\end{itemize}

\subsubsection{Example 4}

Tateoka proposed the following definition, this looks like `sed'
commands.
\begin{verbatim}
    defparenthesis s/ /g
    defdelete //
    defescape

    Manued was proposed by Ikuo Takeuchi who grieved that editing
    manuscripts by s/Electric mail//email; Later, you use the word
    email, please standardize your terms./g was not smoothly done
    because two dimensional information s/od//of/g document is not
    kept in email.
\end{verbatim}

\begin{itemize}
\item Original version :\\ %
  Manued was proposed by Ikuo Takeuchi who grieved that editing
  manuscripts by Electric mail was not smoothly done because two
  dimensional information od document is not kept in email.
\item Revised version : \\ %
  Manued was proposed by Ikuo Takeuchi who grieved that editing
  manuscripts by email was not smoothly done because two dimensional
  information of document is not kept in email.
\end{itemize}

%% Manued Software
%% Copyright (C) By Hitoshi Yamauchi
%% Today, The dummy Professer is not here. No one knows where he
%% is. Even the office does not know.
%% He did not know the word of responsibility. Why such a professor is
%% not be fired?  1998-2-4(Wed)

\newpage
\begin{rawhtml}
<hr>
\end{rawhtml}

\section{A Supporting tool \texttt{manued.el} for Manued.}

\subsection{Download}

You can find the latest version and documents in the following Web sites.

\begin{itemize}
\item \htmladdnormallink{https://github.com/emacsmirror/manued}{https://github.com/emacsmirror/manued}
\end{itemize}


\subsection{Short history of \texttt{manued.el}}

\texttt{manued.el} is an Emacs lisp program for supporting the Manued
proofreading method. Manued method itself is a proofreading method for
correcting electric texts. You can use it even on a paper. But if I had
a program, then I would like to have the following features: emphasizing
manued commands (e.g. using colors), extracting the revised text and the
original text automatically. \texttt{manued.el} offers to add colors to
manued commands, to search correction points, to extract the revised
text and the original text, and to assist to input manued commands.

I have been interested in Lisp language since I took the Masahiko Sato's
class, functional language.  At that time, I have not written any Lisp
programs. I sometimes tried to write a lisp program, but I always failed
because I don't have a killer problem which I really want to solve and I
had no friend who is a lisper.

Several years are passed, I must exchange manuscripts to correct with
research colleagues via email. It was so cumbersome. I found it was so
difficult to revise the manuscript via email, especially just before the
deadline of our paper. First, I wrote my manuscript on my computer, then
I printed it out, checked it by red pen, then, I inputted correction
points back to the computer.

Once I supervised a student. One day he send me his draft of his master
thesis. The time of correcting work reaches several hours and the time
of input work was taken more hours. I used the conventional method for
correcting manuscript. However, there were so much errors and we both
frustrated.

Then, I joined the programming symposium and heard the idea of manued. I
was immediately fascinated this idea. I thought this is the program I
should write and my natural choice is emacs lisp since I am a user of
the editor. Professor Takeuchi commented that ``This program should be
very easy to implement. I think the fourth grade student entered his/her
graduation thesis can be accomplish to write this program.'' I thought
this should be a good exercise to start writing lisp program for
me. Fortunately, Atusi Maeda, a great lisper helped me a lot.


\subsection{Installation}

\texttt{manued.el} is based on Emacs Version 20.7 or later.

\begin{enumerate}
\item Put \texttt{manued.el} the directory which is included
  load-path.
\item Add next line to `.emacs'.
\begin{verbatim}
  (autoload 'manued-minor-mode "manued" "manuscript editing mode" t)
\end{verbatim}
\end{enumerate}

\subsection{How to use \texttt{manued.el}}

Open any text by emacs, then type `M-x manued-minor-mode' to enter the
manued-minor-mode. The following commands are available. Of course you
can customize the key bindings.

\begin{itemize}
\item Starting mode and others
  \begin{itemize}
  \item M-x manued-mode\\ %
    This is alias of next manued-minor-mode, see next item.
  \item M-x manued-minor-mode\\ %
    Enter manued minor mode. \texttt{manued.el} after version 0.9.0
    has only minor mode. The author usually uses the manued mode with
    TeX-mode, LaTeX-mode, indented-text-mode and mail draft mode only.
    Then, for easiness of maintenance for this code, the author
    abolished manued-major-mode.
  \item M-x manued-show-version\\ %
    Show the version of \texttt{manued.el}.
  \item M-x manued-insert-header\\ %
    Insert the def-command header for manued command.
  \end{itemize}

\item Hi-light
      \begin{itemize}
       \item M-x manued-hilit-older\\ %
	     Hi-light and the original part of document in delete
	     command is emphasized.
       \item M-x manued-hilit-newer\\ %
	     Hi-light and the revised part of document in delete command
	     is emphasized.
       \item M-x manued-recenter-hilit\\ %
	     Set current point to the center of the screen and re-hilit
	     the buffer.
      \end{itemize}

\item Extracting document
      \begin{itemize}
       \item M-x manued-show-older-in-manued-buffer\\ %
	     Extract the original document from current buffer to
	     another buffer.
       \item M-x manued-show-newer-in-manued-buffer\\ %
	     Extract the revised document from current buffer to another
	     buffer.
       \item M-x manued-show-older-region\\ %
	     Extract the original document from a designated region of
	     current buffer. The designated region will be changed.
       \item M-x manued-show-newer-region\\ %
	     Extract the revised document from a designated region of
	     current buffer. The designated region will be changed.
       \item M-x manued-eval-last-manuexp\\ %
	     Extract the revised document from the prior manued command.
       \item M-x manued-set-pretty-print-on \\
	     Set pretty print mode on.
       \item M-x manued-set-pretty-print-off \\
	     Set pretty print mode off.
      \end{itemize}

 \item Insert the Manued commands
      \begin{itemize}
       \item M-x manued-insert-swap-command\\ %
	     Insert swap command. The `A' part of swap command is filled
	     with designated region. When
	     \textit{manued-is-swap-command-with-comment-on} is `t',
	     comment symbol is inserted together.

       \item M-x manued-insert-swap-command-toggle-comment\\ %
	     Almost same as \textit{manued-insert-swap-command}, but
	     condition of comment insertion is reversed. Comment symbol
	     is inserted when
	     \textit{manued-is-swap-command-with-comment-on} is nil.

       \item M-x manued-insert-delete-command\\ %
	     Insert delete command. The `delete' part of delete command
	     is filled with designated region. When
	     \textit{manued-is-delete-command-with-comment-on} is `t',
	     comment symbol is inserted together.

       \item M-x manued-insert-delete-command-toggle-comment\\ %
	     Almost same as \textit{manued-insert-delete-command}, but
	     condition of comment insertion is reversed. Comment symbol
	     is inserted when
	     \textit{manued-is-delete-command-with-comment-on} is nil.

       \item M-x manued-insert-comment-region]\\ %
	     Insert comment command. The designated region is surrounded
	     with comment command.

       \item M-x manued-insert-comment-at-point\\ %
	     Insert comment command at the current point.
      \end{itemize}

 \item Search
      \begin{itemize}
       \item M-x manued-next-l-parenthesis\\ %
	     Search next manued command from current point.

       \item M-x manued-previous-l-parenthesis\\ %
	     Search previous manued command from current point.
      \end{itemize}
\end{itemize}

\subsection{Customizing \texttt{manued-minor-mode}}

\textit{Manued.el} provides two main customizing methods. One of them is
to change the emacs lisp variables. I will show the list of customizable
variables later. Another method is to use the hook function.  The entry
function `manued-minor-mode' calls
\begin{verbatim}
        (run-hooks 'manued-minor-mode-hook)
\end{verbatim}
after all initialization is done. So, you can customize
\textit{manued-minor-mode} by using this hook function. (However, now
I did not use this feature, this is only available.)

One of recommended methods is that you put the file
\texttt{.manued.el} in your home directory and set variables in the
file, then you read the file from \texttt{.emacs}. For this, add your
\texttt{.emacs} file to next line.
\begin{verbatim}
        (load-file "~/.manued.el")
\end{verbatim}


\subsection{Customizable variables}

\subsubsection{Manued definition commands of correction symbols (def-command)}

When you set the following variables in \texttt{.emacs} file or set by
any other method, your definition is valid except when def-commands are
appeared in manued document. The default values of def-command is also
shown in `()'.

\begin{itemize}
\item manued-l-parenthesis-str\\ %
  The default value of start symbol of manued command. (default
  `\texttt{[}')
\item manued-r-parenthesis-str\\ %
  The default value of end symbol of manued command. (default
  `\texttt{]}')
\item manued-delete-str\\ %
  The default value of delete symbol of manued command. (default `\texttt{/}')
\item manued-swap-str\\ %
  The default value of swap symbol of manued command. (default `\verb-|-')
\item manued-comment-str\\ %
  The default value of starting comment symbol of manued
  command. (default `\texttt{;}')
\item manued-escape-str\\ %
  The default value of escape symbol of manued command. (default
  `\verb-~-')
\item manued-order-str\\ %
  The default applying order of delete command. (default
  `older-first') The candidate variables are ``older-first'',
  ``newer-last'' or ``older-last'', ``newer-first''. The form of
  delete command of manued is [first/last], you must decide which is
  the revised part (= newer-part). See the item
  \htmlref{deforder}{item:deforder} in section
  \ref{sec:manueddefcommand}.
\end{itemize}

\texttt{manued.el} checks that these def-commands have consistency or
not. If there is no consistency (e.g. two def-commands has same
symbol.), \texttt{manued.el} may warn it. Then you should change
definition of def-commands.


\subsubsection{hi-light colors}
% <label id="HilightColor">

There are three choices for using color. Set following variable to
\texttt{manued-use-color-hilit}.

\begin{itemize}
 \item Use the color as possible as it can.(\texttt{t})
 \item Never use the color.(\texttt{nil})
 \item Following the font-lock-mode status.(\texttt{'follow-font-lock-mode})
\end{itemize}

Example.
\begin{verbatim}
(setq manued-use-color-hilit 'follow-font-lock-mode)
\end{verbatim}

A list which has 6 color names are used to specify the hi-lighting
colors. Emacs teaches us what color names are available by typing `M-x
list-colors-display'. Each elements of the list means as follows.

\begin{verbatim}
'(delete-first-color delete-last-color
  swap-alpha-color swap-beta-color swap-gamma-color
  comment-color command-color)
\end{verbatim}

Where,
\begin{itemize}
\item  delete-first-color, delete-last-color\\ %
  Specify colors of text in delete command as
  \verb|[|delete-first-color/delete-last-color\verb|]|.
\item  swap-alpha-color, swap-beta-color, swap-gamma-color \\ %
  Specify colors of text in swap command as \verb|[|swap-alpha-color |
  swap-beta-color | swap-gamma-color\verb|]|
\item  comment-color \\ %
  Specify the text color in comment region.
\item  command-color\\ %
  Specify the manued command symbol color. As default, manued command
  symbols are `\verb|[|',`/',`;',`\verb|]|'.
\end{itemize}

For example, according to default values, this color list is set with
\begin{verbatim}
 '(red gray60 blue red green4 BlueViolet gray60).
\end{verbatim}
The variables of keeping default color list is next.

\begin{itemize}
\item manued-first-hilit-color-list\\ %
  The color list of when first part is emphasized in delete command.
\item manued-last-hilit-color-list\\ %
  The color list of when last part is emphasized in delete command.
\end{itemize}

\subsubsection{Pretty print}
\label{sec:PrettyPrintCustomize}

\begin{itemize}
 \item manued-pretty-print-format-delete-list\\
       Format list of pretty print for delete command.\\
       (default \verb-'("\\textnormal{%s}" "\\textbf{%s}" "\\textit{%s}")-)

 \item manued-pretty-print-format-swap-list \\
       Format list of pretty print for swap command.\\
       (default
       \verb|'("\\textbf{%s}---" "\\textbf{%s}" "---\\textbf{%s}" "(\\textit{%s})")|

 \item manued-pretty-print-on-p\\
       If \texttt{t}, the document extraction command gives you the
       output with pretty print style. (default \texttt{nil})
       This variable can be changed with commands
       (\verb|M-x manued-set-pretty-print-on|,
       \verb|M-x manued-set-pretty-print-off|).
       If you want to get the output as always pretty print style,
       set this variable to \texttt{t}.

 \item manued-pretty-print-null-comment-out-p\\
       When A manued expression has no comment, pretty print does not
       output the comment part format strings. However, if you want to
       output the comment part format whether a manued expression has or
       doesn't have a comment, you can set this variable \texttt{t}.
       Notice: this variable is not buffer local.

\end{itemize}


\subsubsection{Other customization}

\begin{itemize}
\item manued-mode-syntax-table\\ %
  Default value of syntax-table (default : text-mode-syntax-table)
\item manued-mode-abbrev-table\\ %
  Default value of abbrev-table (default : text-mode-syntax-table)
\item manued-mode-map, manued-minor-mode-map\\ %
  Default value of keymaps
\item manued-is-delete-command-with-comment-on\\ %
  When this value is `t', \textit{manued-insert-delete-command}
  inserts delete command with comment symbol.
\item manued-is-swap-command-with-comment-on\\ %
  When this value is `t', \textit{manued-insert-swap-command} inserts
  swap command with comment symbol.
\item manued-is-auto-insert-header\\ %
  This value is a list of \texttt{'(auto-insert query-when-insert
  insert-point)}.  This value controls inserting def-command and
  behavior of inserting when no def-command is found in manued
  document.
  \begin{itemize}
  \item \texttt{auto-insert} : When this is t��and when entering
    manued-minor-mode, \texttt{manued.el} tries to insert def-command
    header if def-command is not found in manued document.
  \item \texttt{auto-insert} (This is effective iff
    \texttt{auto-insert} is t.) : When this is t, \texttt{manued.el}
    asks insert or not to user with interactively.
  \item \texttt{insert-point} (This is effective iff
    \texttt{auto-insert} is t.) : There are three cases according to
    the value.
    \begin{itemize}
    \item t : Insert header at current cursor position of the buffer.
    \item nil : Insert header at the top of the buffer (point-min).
    \item number : Insert header at this number position.
    \end{itemize}
  \end{itemize}
\item manued-ask-if-formatted-buffer-is\\ %
  When this value is t, and when document extracted buffer is already
  exist, ask user to overwrite the buffer or not.
\item manued-defcommand-head-str-list\\ %
  List of string which can add top of def-command. Default values are
  ``\%'' and ``\%\%'', since you can put def-command into the TeX
  comments.
\end{itemize}

\subsection{Bugs}

\begin{itemize}
 \item Current version (0.9.5) does not use \texttt{font-lock-mode}, so,
       some environment can not be hi-lighted.
\end{itemize}

\subsection{Copying}

\texttt{manued.el} (C) Copyright 1998-2012 Hitoshi Yamauchi

This software is free software. There is absolutely no warranty about
this program. This software can be redistributed only under GNU
copyleft (most recently version of GNU General Public License).

Suggestions, comments and bug reports are welcome to

\newpage
\begin{rawhtml}
<hr>
\end{rawhtml}

\section{Manued Tips}

\begin{itemize}
 \item \htmlref{TeX and Manued (How to compile a TeX file
       which extracted by show newer command immidiately.)}
       {sec:ImmediateShownewerTex}
       \begin{itemize}
	\item \htmlref{When def-command harms TeX compilation.}
	      {sec:harmDefcommand}
	\item \htmlref{When TeX option is deleted with manued processing.}
	      {sec:harmDefaultLRParenthes}
       \end{itemize}
\end{itemize}


\label{sec:ImmediateShownewerTex}
\subsection{TeX and Manued}


How to compile TeX file which is extracted by show-newer-command
immediately.

\label{sec:harmDefcommand}
\subsubsection{When def-command harms TeX compilation.}

I would like to latex compile the output of show-newer-command immediately.
(First, you apply manued-show-newer-in-manued-buffer to the current
buffer. Second, you write *manued* buffer which is result of the command
to foo.tex file by using M-x write-file. Finally you compile `foo.tex'
file.) However the compilation may fail because of manued
def-commands. One solution is not to use def-command in the manuscript
and to set manued-def-* variables by elisp startup
file(\textit{.emacs}).  But one of the meanings of def-command is to
keep consistency between the def-command symbols and the manued commands
within the document. Therefore, it is not a good idea to eliminate
def-command header, especially a manuscript is exchanged via email since
this is just the case that such consistency should be kept.

\texttt{manued.el} can recognize extended def-command which is
attached with some strings at top of def-command. You can set the head
strings of def-command to \textit{manued-defcommnd-head-str-list}. By
default, this list includes `\%' and `\%\%' which is the comment symbols
of TeX. Next code is recognized by \texttt{manued.el} as
def-command. However, a manued command must be put at the beginning of
line and should not include blank between the head string with the
manued command. You can also customize
\textit{manued-defcommnd-head-str-list} to change the header string.

\begin{verbatim}
        %%defLparenthesis [
        %%defRparenthesis ]
\end{verbatim}

\label{sec:harmDefaultLRParenthes}
\subsubsection{When TeX option is deleted with manued processing}

Because the default manued command starts at `\verb|[|' and ends with
`\verb|]|', this is same as the TeX option argument. I always use
Japanese environment, I change the `\verb|[|' and `\verb|]|' to 2bytes
parentheses character. However, people use ASCII environment may not
use such method. So, one of this solution uses some special strings
instead of `\verb|[|' and `\verb|]|'. Next setting is an example.

\begin{verbatim}
        (setq manued-l-parenthesis-str ":[:")
        (setq manued-r-parenthesis-str ":]:")
\end{verbatim}

However this seems not so good solution. I actually use a multi-byte
characters, but that also depends on your environment.

\newpage
\begin{rawhtml}
<hr>
\end{rawhtml}

\section{Manued Links}

\subsection{Other implementation of manued.el (By Prof. WADA)}

Takeuchi and I introduce another implementation of manued on ``bit''
No. 12, 1999. Prof. Wada also implemented it on emacs.

This section introduces the implementation of manued by Prof. Wada.
This implementation focuses to ``see'' the manued document and much
smaller than my implementation. This compactness allows to send a
manuscript with the code to whom a person did not install
\texttt{manued.el} and the person can instantly use it.

\subsubsection{Implementation of manued by Prof. Wada.}
\begin{verbatim}
;;
;; manued : editing manuscripts for Emacs
;; 1998 Eiiti Wada
;; manued.el, Copyright (C) 1999 Eiiti Wada��
;;
;; (YAMAUCHI, Hitoshi translates the comments to English for explanation.)
;;
;;[A\B]   means `substitute A with B.'
;;[A\]    means `delete A.'
;;[\B]    means `insert B.'
;;[A|B|C] means `reorder ABC to CBA.'
;;[A||B]  means `reorder AB  to BA.'
;;[@C]    means comment.
;;
;; Someone said that it's hard to read the manued document, so I quick
;; hacked a tool for easy browsing. Let's try it. This tool works with
;; X Window version of Mule editor.
;;
;; First, put this emacs lisp program in the mule buffer.
;; Second, type `esc' `esc', then you will see the prompt `Eval:'
;; at the minibuffer, then type (eval-current-buffer) to evaluate whole
;; the program.
;;
;; Next, put the text of your email in the mule buffer, type `esc' `esc'
;; and (manued). Then, you can see the colored A, B, C potion in the
;; buffer.
;;
;; If you don't like these colors, type `esc' `esc' again and execute
;; (list-faces-display), then you can get the name list of enable
;; colors.  According to the color list, change the color definition
;; part --- that is in the head part of this program --- to your
;; favorite colors.
;;

(make-face 'alpha)
(set-face-foreground 'alpha "blue")
(make-face 'beta)
(set-face-foreground 'beta "red")
(make-face 'gamma)
(set-face-foreground 'gamma "green4")
(setq facealpha 'alpha)
(setq facebeta  'beta)
(setq facegamma 'gamma)

(defun proc-range (beg end)
 (let* (p0 p1 p2 ch)
       (goto-char beg)
       (setq beg (search-forward "[" end t))
       (while beg (proc-sec) (setq beg (search-forward "[" end 0)))))

(defun proc-sec nil
 (setq p0 (re-search-forward "[]|\[{]" end t))
 (setq ch (string-to-char (buffer-substring (1- p0) p0)))
 (cond ((= ch ?[) (skip) (proc-sec))
       ((= ch ?\\) (findrsb) (alpha beg p0) (beta p0 p2))
       ((= ch ?|) (findvl) (findrsb) (alpha beg p0) (gamma p0 p1) (beta p1 p2))
       ((= ch ?]) (gamma beg p0))
       ((= ch ?{) (skipcb) (proc-sec))))

(defun skip nil
 (let* (p ch)
       (setq p (re-search-forward "[][{]" end t))
       (while p
              (setq ch (string-to-char (buffer-substring (1- p) p)))
              (cond ((= ch ?]) (setq p nil))
                    ((= ch ?[)
                     (skip)
                     (setq p (re-search-forward "[][]" end t)))
                    ((= ch ?{) (skipcb))   ))))
(defun skipcb nil
 (let* (p ch)
       (setq p (re-search-forward "[}{]" end t))
       (while p
              (setq ch (string-to-char (buffer-substring (1- p) p)))
              (cond ((= ch ?}) (setq p nil))
                    ((= ch ?{)
                     (skipcb)
                     (setq p (re-search-forward "[}{]" end t)))))))

(defun findrsb nil
 (setq p2 (re-search-forward "[][{]" end t))
 (setq ch (string-to-char (buffer-substring (1- p2) p2)))
 (cond ((= ch ?[) (skip) (findrsb))
       ((= ch ?]))
       ((= ch ?{) (skipcb) (findrsb))))

(defun findvl nil
 (setq p1 (re-search-forward "[|[{]" end t))
 (setq ch (string-to-char (buffer-substring (1- p1) p1)))
 (cond ((= ch ?[) (skip) (findvl))
       ((= ch ?|))
       ((= ch ?{) (skipcb) (findvl))))

(defun alpha (l r)
 (put-text-property l (1- r) 'face facealpha)
 (proc-range l (1- r)))

(defun beta  (l r)
 (put-text-property l (1- r) 'face facebeta)
 (proc-range l (1- r)))

(defun gamma (l r)
 (put-text-property l (1- r) 'face facegamma)
 (proc-range l (1- r)))

(defun manued nil (proc-range 1 (point-max)))
\end{verbatim}

I thank Prof. Wada Eiiti for permitting to public this code.

\subsection{Other implementations}

There is an implementation for a text editor named ``Hidemaru'' on
Windows 95 and later.

\subsection{Manueder (By Suzuki)}

\htmladdnormallink{Manuediff}
{http://hibiki.miyagi-ct.ac.jp/~suzuki/comp/export/manuediff.html}
is a diff program which outputs manued format. Example:
\begin{verbatim}
  manuediff file1 file2 > diff.with.manued.format
\end{verbatim}

\newpage
\begin{rawhtml}
<hr>
\end{rawhtml}


\section{Acknowledgements}

Thanks to Takeuchi Ikuo to offer ``Manuscript Editing on E-mail'' and
its electric text.

I also thanks to Eiiti Wada who gave me an interesting surprise by
the concise implementation.

I also should say ``thank you'' to Atusi Maeda who teaches me the lisp
language. I had no knowledge about lisp when I started to implement
this.

Some people gave fruitful suggestions, comments with working codes to
\texttt{manued.el} implementation. I will thank for all such people,
especially the following people.  (Alphabetical order)
\begin{itemize}
 \item Ken-ich Suzuki (manuediff)
 \item Mikio Nakajima (bug fix etc. for version 0.9.3)
 \item Stefan Monnier  (comments for version 0.9.1)
 \item Takao Kawamura (bug fix etc. for version 0.9.3)
\end{itemize}

\nocite{bib:Takeuchi1986}
\nocite{bib:Aoyagi1997}
\nocite{bib:Hirose1999}
\nocite{bib:elispinfo}
\addcontentsline{toc}{section}{References}
\bibliographystyle{plain}
\bibliography{reference-e}
\printindex

\begin{rawhtml}
<address>
Copyright (C) 1998-2012 yamauchi, Hitoshi<br>
Most recent update : <!--#echo var="LAST_MODIFIED" --> :
</address>
\end{rawhtml}

\end{document}


%%% Local Variables:
%%% mode: yatex
%%% TeX-master: t
%%% End:
